\begin{exercice*}
    Pour chaque question, entoure la (ou les) bonne(s) réponse(s).

    \begin{enumerate}
        \item Un jeu de 32 cartes comporte quatre rois. On tire une carte au hasard. Quelle est la probabilité d'obtenir un roi ?

        \[\dfrac{1}{8} \qquad \dfrac{1}{32} \qquad \dfrac{4}{32} \qquad \dfrac{4}{28}\]

        \item Une urne contient deux boules rouges, cinq boules vertes et trois boules bleues. Quelle est la probabilité d'obtenir une boule rouge ?

                  \[\dfrac{2}{8} \qquad \dfrac{2}{10} \qquad \dfrac{8}{2} \qquad \dfrac{1}{5}\]

        \item Une urne contient cinq boules bleues numérotées de 1 à 5, trois boules blanches numérotées de 1 à 3 et deux boules noires numérotées de 1 à 2. Quelle est la probabilité d'obtenir un 3 ?

                  \[\dfrac{3}{10} \qquad \dfrac{2}{10} \qquad \dfrac{2}{5} \qquad \dfrac{1}{5}\]
    \end{enumerate}
\end{exercice*}
\begin{corrige}
    %\setcounter{partie}{0} % Pour s'assurer que le compteur de \partie est à zéro dans les corrigés
    %\phantom{rrr}    
    \begin{enumerate}
        \item La probabilité d'obtenir un roi est de $\dfrac{4}{32}$, car il y a 4 rois sur les 32 cartes au total.
        \item La probabilité d'obtenir une boule rouge est de $\dfrac{2}{10}$, car il y a 2 boules rouges sur les 10 boules au total (2 rouges + 5 vertes + 3 bleues).
        \item La probabilité d'obtenir un 3 est de $\dfrac{3}{10}$, car il y a 3 boules numérotées 3 (une bleue, une blanche et une noire) sur les 10 boules au total (5 bleues + 3 blanches + 2 noires).
    \end{enumerate}
\end{corrige}

